\documentclass[11pt]{article}

\usepackage[margin=1in]{geometry}
\usepackage{array}
\usepackage{booktabs}
\usepackage{longtable}
\usepackage{tabularx}
\usepackage{xcolor}
\usepackage{hyperref}

\hypersetup{
  colorlinks=true,
  linkcolor=black,
  urlcolor=blue
}

\setlength{\parindent}{0pt}
\setlength{\parskip}{6pt}
\renewcommand{\arraystretch}{1.18}

\newcolumntype{Y}{>{\raggedright\arraybackslash}X}
\newcommand{\markcell}[2]{\textbf{#1/#2}}

\begin{document}

\begin{center}
  {\Large Geography Paper 2 TZ1 HLSL 2024: Student Work Marking Report}\\[4pt]
  {\normalsize Marked against \texttt{Geography\_paper\_2\_\_TZ1\_HLSL\_markscheme-2024.pdf}}
\end{center}

\section*{Overall Result}

\begin{tabularx}{\textwidth}{@{}p{0.34\textwidth} c c Y@{}}
\toprule
Section & Marks awarded & Maximum & Comment \\
\midrule
Question 1: Changing population & 9 & 10 & Accurate resource-reading answers and thoughtful ageing-policy explanations; one political displacement answer needed an explicit link to people fleeing or relocating. \\
Question 2: Global climate & 9 & 10 & Strong distribution and technology answers; the permafrost methane answer needed the release mechanism to be stated. \\
Question 3: Resource consumption and security & 10 & 10 & Clear definitions, developed energy-availability explanations, and a well-applied Kiribati water-security answer. \\
Question 4: Sustainable fashion infographic & 9 & 10 & Excellent six-mark appraisal; the short graph answer missed the initial fall / changing rate in the data. \\
Question 5: Megacity growth essay & 7 & 10 & Relevant, structured and comparative, but too general and lightly evidenced for the top band. \\
\midrule
\textbf{Total} & \textbf{44} & \textbf{50} & Main targets: make every causal pathway explicit, describe graph changes in stages, and strengthen Section C with more precise evidence and comparison. \\
\bottomrule
\end{tabularx}

\section*{Question-by-Question Marks}

\begin{longtable}{@{}p{0.12\textwidth} p{0.12\textwidth} p{0.31\textwidth} p{0.37\textwidth}@{}}
\toprule
Question & Mark & Awarded for & Feedback \\
\midrule
\endfirsthead
\toprule
Question & Mark & Awarded for & Feedback \\
\midrule
\endhead
1(a)(i) & \markcell{1}{1} & Gives 30\%. & Correct; within the accepted range. \\
1(a)(ii) & \markcell{1}{1} & Identifies prevent abuse. & Correct policy. \\
1(b) Reason 1 & \markcell{2}{2} & Explains that the health of the ageing population affects the level of healthcare investment and intervention. & Valid reason for policy variation, developed through different healthcare needs and costs. \\
1(b) Reason 2 & \markcell{2}{2} & Explains that the physical demands of employment affect whether later retirement or work incentives are viable. & Valid and developed reason for different ageing-society policies. \\
1(c) Environmental & \markcell{2}{2} & Identifies floods/natural disasters and links them to eroding soil and weakening building foundations. & This is enough for environmental displacement through damage to homes. \\
1(c) Political & \markcell{1}{2} & Identifies conflict caused by political polarization. & Add the second step: people flee to safer areas or are forced to relocate internally. \\
\midrule
2(a) & \markcell{2}{2} & Describes a northward shift and a larger/more northern distribution by 2100 compared with the central band in 1990. & Clear spatial comparison. \\
2(b) Reason 1 & \markcell{1}{2} & Identifies melting permafrost as global warming continues. & To secure the second mark, state that methane stored in permafrost is released when it thaws. \\
2(b) Reason 2 & \markcell{2}{2} & Explains that increased beef consumption increases cow numbers, and cows produce methane during digestion. & Full cause-effect chain. \\
2(c) Way 1 & \markcell{2}{2} & Explains that electric vehicles reduce combustion of fossil fuels and reduce carbon emissions. & Valid mitigation technology. \\
2(c) Way 2 & \markcell{2}{2} & Explains that in vitro meat reduces the need for livestock and therefore reduces methane emissions. & Valid technological approach to reducing climate-change effects through lower emissions. \\
\midrule
3(a) & \markcell{2}{2} & Defines resource stewardship as sustainable resource use so future population demands can be met. & Complete definition. \\
3(b) Way 1 & \markcell{2}{2} & Explains that economic development can fund costly equipment needed to use renewable energy. & Valid link between development, investment and energy availability. \\
3(b) Way 2 & \markcell{2}{2} & Explains that economic development improves access to trade, allowing regions to buy energy from energy-rich places. & Valid import/trade route to increased energy availability. \\
3(c) Factor 1 & \markcell{2}{2} & Uses Kiribati and explains flooding/saltwater contamination of freshwater reserves and wells. & Well-applied place-specific water-security factor. \\
3(c) Factor 2 & \markcell{2}{2} & Explains that overcrowding on the main atoll makes it harder to provide enough potable water. & Valid demand-side factor, relevant to Kiribati. \\
\midrule
4(a) & \markcell{1}{2} & Gives the overall increase with approximate values from 2014 and 2020. & For full marks, describe the pattern in stages: the initial decrease, the later increase, and/or the slower rate after 2018. \\
4(b) & \markcell{2}{2} & Suggests that the selected cities may not represent views in other cities. & Valid bias with development. \\
4(c) & \markcell{6}{6} & Uses evidence for positive impacts, including reduced carbon, water and waste footprints, recycling/reuse, landfill/incineration reduction and circular economy benefits; also weighs the GDP and employment downside of disposable fashion. & Strong balanced appraisal with a clear final judgement. Full marks. \\
\midrule
5 & \markcell{7}{10} & Addresses individuals, governments/society and the environment; includes positive and negative consequences, a Jakarta slum example, and a relevant overall judgement. & This reaches the 7--8 band because it compares who feels the impacts. It stays at the lower end because the evidence is limited, statistics are absent, and the answer needs more precise comparison of individuals, wider society, rural areas, wealth groups and environmental systems. \\
\bottomrule
\end{longtable}

\section*{Teaching Feedback}

\textbf{What is working well}
\begin{itemize}
  \item Short-answer command terms are handled confidently, especially definitions and two-mark explanation structures.
  \item The student usually builds clear cause-effect chains rather than listing isolated points.
  \item Place application is strong in Question 3(c): Kiribati is used accurately for water-security pressures.
  \item The Question 4(c) appraisal is balanced and uses evidence from both sides of the infographic.
  \item The Section C essay is organized into a comparative structure, which is the right direction for a ``to what extent'' question.
\end{itemize}

\textbf{Main improvement targets}
\begin{itemize}
  \item Make the final causal step explicit. For displacement, state that people flee, evacuate, or are forced to relocate; for permafrost, state that stored methane is released.
  \item For graph-description questions, describe the pattern in stages rather than only giving the start and end values.
  \item In six- and ten-mark answers, keep using evidence, but add more precise statistics, named examples and comparison.
  \item In Section C, compare individuals with society/government more directly, and consider differences by wealth group, place, rural/urban impacts and time scale.
\end{itemize}

\section*{Suggested Next Practice}

\begin{itemize}
  \item Rewrite Question 1(c)'s political factor in two sentences: factor, then how it forces internal movement.
  \item Rewrite Question 2(b) Reason 1 using the full chain: warming, thawing permafrost, methane release.
  \item Redo Question 4(a) with at least two separate changes from the graph, including the small fall from 2014 to 2015.
  \item Re-plan Question 5 with a comparison grid for individuals, society/government and the environment. Add at least one statistic or named place detail to each column.
\end{itemize}

\section*{Model Improvements}

\textbf{Question 1(c) stronger political factor:}
Conflict between political groups can make neighbourhoods unsafe. People may then move internally to safer regions of the same country to avoid violence or persecution.

\textbf{Question 2(b) stronger permafrost chain:}
Rising temperatures thaw Arctic permafrost. Methane stored in the frozen ground is then released into the atmosphere, increasing methane emissions.

\textbf{Question 4(a) stronger graph description:}
Sustainable clothing sales revenue fell slightly from 2014 to 2015, then increased overall to about 58 million GBP by 2020. The rate of increase was slower from 2018 to 2020 than in the earlier growth period.

\textbf{Question 5 stronger essay route:}
Megacity growth is strongly felt by individuals through housing, health, employment, transport and inequality, but wider society and governments also face infrastructure costs, pollution, service pressure and economic restructuring. A higher-band answer should weigh these scales of impact using named evidence, then judge whether individual impacts are more direct, more severe, or simply more visible than societal impacts.

\section*{Notes}

\begin{itemize}
  \item The marked student-work source is the 2024 TZ1 HLSL image set, \texttt{student-work-1.jpg} through \texttt{student-work-8.jpg}. These images show answers for Section A Questions 1--3, Section B Question 4, and Section C Question 5. Question 6 was not attempted and was not marked.
  \item This report records marking judgements and teaching feedback only; it does not reproduce the exam paper or markscheme.
\end{itemize}

\end{document}
